\section{Discussion and Limitations}
\label{sec:discussion}

\subsection{What is established, and what scales with data}
The layered validation of \Cref{sec:validation} separates two kinds of claims, and the
distinction should govern how the results are read. What the evidence establishes: the
transforms are dimensionally guarded and fail loudly rather than silently; every
implemented forward mapping is accepted by the economy solve; the response directions
follow the stated mechanisms (a resource cost contracts, the factor-share lever is not
a capital draw-in, the survival gain skews old and therefore carries almost no output);
and the composition of the coupled response is pinned by construction and confirmed by
matched runs. What remains conditional is every \emph{level}. The unit/deflator bridge
between the energy model's money and the economy's num\'eraire is not yet audited, so
any currency magnitude that crosses it --- carbon revenue, investment cost --- is
illustrative. The electricity-cost signal is ratio-based and thus immune to the
deflator, but it is a levelized cost, not a marginal price. The morbidity multiplier is
a placeholder pending a country dose--response estimate. This is why every magnitude in
\Cref{sec:results} travels as a relative change with its evidence base named.

\subsection{The transmission choice is first-order}
The single most consequential specification in the system is not a parameter but a
routing. The composite's inter-industry leg treats the sectoral electricity-cost change
as an exogenous unit-cost shift, scaled by input--output exposure; it is a reduced-form
mapping, not a factor-demand system derived from the production block. Nearly the whole
coupled response arrives through it, and a structural alternative that routes the same
cost signal through sectoral productivity reverses the aggregate sign
(\Cref{sec:decomposition}). The two routings embody different economic claims about
what a costlier power system does to the rest of the economy, and the present framework
cannot adjudicate between them because electricity is not a priced input inside
production. Until it is, the choice is an assumption --- which is why the headline is
never quoted without the transmission named.

\subsection{One pass, not yet a fixed point}
The emitted inputs --- the carbon-price artifact, the planning rate, the revised demand
path --- bind nothing until a second energy-model solve consumes them. The mechanical
seam for that second solve is built and tested: a controlled demand patch survives
copy, re-solve, and read-back, deterministically. What does not yet exist is the outer
controller --- damping, a convergence criterion, and control of propagated error across
repeated solutions --- and the validation table lists closed-loop iteration as untested
for exactly that reason. Until it lands, the system is a forward evaluator: the
reported responses are conditional on an energy plan that has not yet been allowed to
react to them.

The cost of closing the loop is expected to be bookkeeping rather than compute, for
one reason on each side. Each iteration re-solves the energy model's linear program
with only a handful of coefficients changed --- the emitted demand path, discount
rate, and carbon penalty --- and a solver warm-started from the prior optimal basis
reoptimizes such a neighbouring problem at what is typically a small fraction of a
cold solve. At the outer level, the exchanged vector is small and smooth, and driving
it to consistency is a fixed-point problem of exactly the kind Anderson mixing
accelerates --- the same scheme the economy model's equilibrium solver already applies
internally, used one level up. Neither technique is exotic, and neither has been
exercised here; the convergence controller above is where both would live.

\subsection{The structural gap behind the proxy}
The economy's production block takes capital and labour; electricity is not a priced
intermediate. Two consequences follow. First, firms cannot substitute away from energy
when it becomes dear, so the cost-push mapping, which imports input--output exposures
from the social accounting matrix, has no behavioural margin behind it. Second,
inter-industry energy flows cannot reallocate endogenously; the exposure weights are
data, not decisions. The proxy is nonetheless a defensible interim: the exposures are
country-sourced, the mapping is transparent enough to audit line by line, and the
structural alternative --- adding a priced energy factor to the host model's
production block --- would modify the host rather than link to it, breaking the
non-invasive constraint the architecture is built on (\Cref{sec:framework}). The
structural ideal belongs on the roadmap, not in the link.

\subsection{Health is a sourced bridge, not an endogenous channel}
The health mapping's inputs are sourced --- the burden total and both age profiles come
from the Global Burden of Disease --- but the translation from power-sector emissions
to that burden is a constructed dose--response scalar whose linearity is assumed, and
the morbidity productivity multiplier is a placeholder. One further point is easy to
miss: because ambient-particulate mortality is concentrated at post-working ages,
aggregate output is nearly invariant to the mortality gain, and GDP is therefore the
wrong summary statistic for this channel. The framework reports lives and returned
working time as first-class outputs beside the aggregates for that reason, not as
decoration.

\subsection{Calibration discipline, checked externally}
Every energy-side parameter the application touched was derived from a named source
rather than tuned toward an observed outcome (\Cref{sec:calibration}), and each
derivation is recorded with its source in the calibration register. That discipline
recently acquired an external check: a subsequent, independently prepared release of
the underlying country energy dataset adopted the same values for the coal capital
cost, the small-modular-reactor capital cost, the large-hydro lifetime, the LNG
chronology, the onshore-wind resource ceiling, and the two-sided land-endowment pin.
Independent convergence on the same sourced values is the strongest evidence available
that the register reflects the energy system being modelled rather than the
calibrator's target.

\subsection{Threats to validity}
Compactly, what would change the conclusions:
\begin{itemize}
\item the transmission assumption --- sign-critical, as shown, and unresolved until
  energy is priced in production;
\item the unaudited deflator --- level-critical for carbon and investment;
\item one-pass conditionality --- the energy plan has not reacted to the economy;
\item the degenerate commodity dual --- the price signal is an average cost, and no
  marginal interpretation survives it;
\item the reduced-form health translation --- magnitude-critical for the health
  channel's non-demographic content;
\item scope --- one country and one scenario pair.
\end{itemize}
